\documentclass[11pt,a4paper]{article}
\usepackage[utf8]{inputenc}
\usepackage{amsmath,amssymb,amsthm}
\usepackage{algorithm}
\usepackage{algpseudocode}
\usepackage{booktabs}
\usepackage{hyperref}
\usepackage{geometry}
\geometry{margin=1in}

\title{RoIAlign and Mask R-CNN: Non-Quantized Bilinear Feature Extraction and Decoupled Instance Segmentation}
\author{Team Scaibu \\ \small Email: \texttt{jaydeep.9664920749@gmail.com} \\ \small Architecture and Core Engine Team}
\date{\today}

\newtheorem{theorem}{Theorem}
\newtheorem{lemma}[theorem]{Lemma}
\newtheorem{proposition}[theorem]{Proposition}
\theoremstyle{definition}
\newtheorem{definition}{Definition}
\newtheorem{remark}{Remark}

\begin{document}

\maketitle

\begin{abstract}
Standard region-of-interest pooling operations (RoIPool) introduce severe quantization artifacts by rounding floating-point region proposals and spatial bin boundaries to discrete integer grid coordinates. This spatial discretization error cascades through network layers, crippling pixel-level instance segmentation accuracy. RoIAlign resolves this fundamental misalignment by entirely eliminating coordinate quantization, sampling exact continuous sub-pixel coordinates through bilinear interpolation. Mask R-CNN extends this aligned representation with a decoupled fully convolutional mask prediction branch. This document formalizes continuous RoI spatial sampling, proves alignment error bounds, presents algorithmic pseudocode, and walks through a verified numerical calculation.
\end{abstract}

\section{Notation and Mathematical Preliminaries}

Let $F \in \mathbb{R}^{C \times H \times W}$ denote a feature map generated with spatial stride $s$ (e.g. $s = 16$), and let an input RoI bounding box be specified in continuous image coordinates as $(x_1, y_1, x_2, y_2) \in \mathbb{R}^4$.

\begin{table}[h]
\centering
\caption{Mathematical Notation for RoIAlign and Mask R-CNN}
\begin{tabular}{lll}
\toprule
\textbf{Symbol} & \textbf{Domain} & \textbf{Semantic Description} \\
\midrule
$k_h \times k_w$ & $\mathbb{N}^2$ & Target pooled spatial grid resolution (e.g. $7 \times 7$ or $14 \times 14$) \\
$(r_x, r_y, r_w, r_h)$ & $\mathbb{R}^4$ & Unquantized continuous RoI box on feature map ($r_x = x_1/s$, etc.) \\
$\Delta_h, \Delta_w$ & $\mathbb{R}_{> 0}^2$ & Continuous bin dimensions: $\Delta_h = r_h / k_h, \Delta_w = r_w / k_w$ \\
$N_s$ & $\mathbb{N}$ & Number of regular sampling points per bin (typically $2 \times 2 = 4$) \\
$(p_y^{(u)}, p_x^{(v)})$ & $\mathbb{R}^2$ & Exact continuous coordinates of the $(u, v)$-th sampling sub-point \\
$y(c, i, j)$ & $\mathbb{R}$ & Extracted feature activation for channel $c$ at pooled bin $(i, j)$ \\
$m \times m$ & $\mathbb{N}^2$ & Predicted instance mask resolution ($28 \times 28$ typical) \\
\bottomrule
\end{tabular}
\end{table}

\section{Problem Definition and Core Concepts}

\subsection{3-Level Explanation}
\begin{enumerate}
    \item \textbf{Intuitive (High School level):} When cutting out a sticker from a magazine, RoIPool rounded your cut to the nearest coarse grid line, tearing the object's edges. RoIAlign uses a laser cutter at the exact sub-millimeter position, blending the four closest pixels perfectly so edges remain sharp and intact.
    \item \textbf{Engineering Level:} RoIPool applied $\lfloor x/16 \rfloor$ twice, resulting in up to 1-pixel feature map misalignment (equal to 16 image pixels). RoIAlign preserves continuous float coordinates $x/16$, divides each bin into regular sample points, and computes values via bilinear interpolation without quantization.
    \item \textbf{Mathematical Level:} For continuous RoI $[x_1/s, x_2/s] \times [y_1/s, y_2/s]$ divided into $k_h \times k_w$ bins, the $(i, j)$-th pooled activation is:
    \begin{equation}
    y(c, i, j) = \frac{1}{N_s} \sum_{u=1}^{\sqrt{N_s}} \sum_{v=1}^{\sqrt{N_s}} \mathcal{B}\left( F_c, y_{\text{bin}}(i) + \left(u - \frac{1}{2}\right)\frac{\Delta_h}{\sqrt{N_s}}, x_{\text{bin}}(j) + \left(v - \frac{1}{2}\right)\frac{\Delta_w}{\sqrt{N_s}} \right)
    \end{equation}
    where $\mathcal{B}(F, y, x)$ represents 2D bilinear interpolation.
\end{enumerate}

\subsection{5-Step Algorithmic Framework}
\begin{enumerate}
    \item \textbf{Continuous Scale Projection:} Map image RoI coordinates $(x_1, y_1, x_2, y_2)$ to feature space by dividing by stride $s$ without rounding.
    \item \textbf{Uniform Continuous Bin Partitioning:} Divide feature RoI into $k_h \times k_w$ equal sub-regions of continuous width $\Delta_w = r_w / k_w$ and height $\Delta_h = r_h / k_h$.
    \item \textbf{Sub-Point Grid Generation:} Place $N_s$ sampling locations inside each bin with coordinates $(y_i + (u - 0.5)\frac{\Delta_h}{\sqrt{N_s}}, x_j + (v - 0.5)\frac{\Delta_w}{\sqrt{N_s}})$.
    \item \textbf{Bilinear Neighbor Interpolation:} Interpolate feature values for all $N_s$ points from 4 surrounding integer grid points.
    \item \textbf{Sample Averaging / Max-Pooling:} Average the $N_s$ sample values to produce the final pooled bin activation $y(c, i, j)$.
\end{enumerate}

\section{Algorithm Specification}

\begin{algorithm}[H]
\caption{RoIAlign Feature Pooling}
\label{alg:roialign}
\begin{algorithmic}[1]
\Require Feature map $F \in \mathbb{R}^{C \times H \times W}$, RoI box $(x_1, y_1, x_2, y_2)$, stride $s$, pooled size $(k_h, k_w)$, samples per bin axis $n_s = 2$.
\Ensure Pooled feature tensor $Y \in \mathbb{R}^{C \times k_h \times k_w}$.
\State $r_{x1} \gets x_1 / s, \quad r_{y1} \gets y_1 / s, \quad r_{x2} \gets x_2 / s, \quad r_{y2} \gets y_2 / s$
\State $r_w \gets \max(r_{x2} - r_{x1}, 1.0), \quad r_h \gets \max(r_{y2} - r_{y1}, 1.0)$
\State $\Delta_w \gets r_w / k_w, \quad \Delta_h \gets r_h / k_h$
\For{$c = 1$ \textbf{to} $C$}
    \For{$i = 0$ \textbf{to} $k_h - 1$}
        \For{$j = 0$ \textbf{to} $k_w - 1$}
            \State $\text{sum} \gets 0.0$
            \For{$u = 0$ \textbf{to} $n_s - 1$}
                \State $p_y \gets r_{y1} + i \cdot \Delta_h + (u + 0.5) \cdot (\Delta_h / n_s)$
                \For{$v = 0$ \textbf{to} $n_s - 1$}
                    \State $p_x \gets r_{x1} + j \cdot \Delta_w + (v + 0.5) \cdot (\Delta_w / n_s)$
                    \State $\text{val} \gets \text{BilinearSample}(F[c], p_y, p_x)$
                    \State $\text{sum} \gets \text{sum} + \text{val}$
                \EndFor
            \EndFor
            \State $Y[c, i, j] \gets \text{sum} / (n_s \cdot n_s)$
        \EndFor
    \EndFor
\EndFor
\State \Return $Y$
\end{algorithmic}
\end{algorithm}

\section{Theoretical Guarantees and Quantization Error Analysis}

\begin{theorem}[Elimination of Spatial Discretization Error]
\label{thm:roialign_error}
Let $\epsilon_{\text{pool}}$ denote the spatial coordinate offset error between continuous image coordinates and feature sampling points. Under standard RoIPool with stride $s$ and $k$ bins:
\begin{equation}
\mathbb{E}[\|\epsilon_{\text{RoIPool}}\|_2] = \mathcal{O}(s)
\end{equation}
whereas under RoIAlign:
\begin{equation}
\|\epsilon_{\text{RoIAlign}}\|_2 \equiv 0
\end{equation}
guaranteeing exact pixel-to-feature spatial registration.
\end{theorem}

\begin{proof}
RoIPool rounds feature box coordinates: $\tilde{x} = \lfloor x / s \rfloor$ and bin sizes $\tilde{\Delta} = \lfloor \tilde{w} / k \rfloor$. In the worst case, round-off accumulates error $\Delta x = |x/s - \lfloor x/s \rfloor| s \in [0, s)$. In contrast, RoIAlign performs continuous affine coordinate mapping $\mathbf{p} = \mathbf{A} \mathbf{x} + \mathbf{b}$ where $\mathbf{A}, \mathbf{b} \in \mathbb{R}^2$, without rounding operators $\lfloor \cdot \rfloor$, preserving exact continuous spatial coordinates.
\end{proof}

\section{Complexity Analysis}

\begin{itemize}
    \item \textbf{Time Complexity:} $\mathcal{O}(C \cdot k_h \cdot k_w \cdot N_s)$ operations per RoI proposal ($N_s = 4$ constant factor).
    \item \textbf{Space Complexity:} $\mathcal{O}(C \cdot k_h \cdot k_w)$ memory per pooled proposal tensor.
\end{itemize}

\section{Exactness and Error Analysis}

The backward pass of RoIAlign distributes gradient $\frac{\partial \mathcal{L}}{\partial Y(c, i, j)}$ back to the 4 continuous neighboring feature pixels through bilinear distance coefficients $G(q, p) = \max(0, 1-|q_x-p_x|)\max(0, 1-|q_y-p_y|)$, maintaining continuous differentiability.

\section{Worked Numerical Example}

Let feature channel $F = \begin{bmatrix} 10.0 & 20.0 \\ 30.0 & 40.0 \end{bmatrix}$ ($H=2, W=2$, stride $s=1$).
RoI box: $(x_1, y_1, x_2, y_2) = (0.0, 0.0, 1.0, 1.0)$, target pooled size $1 \times 1$, $n_s = 1$ (1 central sample).
\begin{itemize}
    \item $r_w = 1.0, r_h = 1.0 \implies \Delta_w = 1.0, \Delta_h = 1.0$.
    \item Single sample point at $u=0, v=0$:
    $p_y = 0.0 + 0.5(1.0) = 0.5$, $p_x = 0.0 + 0.5(1.0) = 0.5$.
    \item Bilinear weights to $(0,0), (0,1), (1,0), (1,1)$ are all $0.25$.
    \item $y(0,0) = 0.25(10) + 0.25(20) + 0.25(30) + 0.25(40) = 25.0$.
\end{itemize}

\section{Numerical Considerations and Edge Cases}

\begin{itemize}
    \item \textbf{Sampling Point Density:} Using $N_s = 4$ ($2 \times 2$) yields practically identical accuracy to dense sampling ($N_s = 16$), while consuming $4\times$ less compute.
    \item \textbf{Mask Branch Loss Isolation:} The Mask R-CNN mask loss $\mathcal{L}_{\text{mask}}$ is evaluated solely on the ground-truth class slice, preventing competition across classes.
\end{itemize}

\section{References}
\begin{enumerate}
    \item He, K., Gkioxari, G., Doll{\'a}r, P., & Girshick, R. (2017). Mask R-CNN. \emph{ICCV}, 2961--2969.
    \item Dai, J., He, K., & Sun, J. (2016). R-FCN: Object detection via region-based fully convolutional networks. \emph{NeurIPS}, 379--387.
\end{enumerate}

\end{document}
